% preamble.tex — 《從高中數學到費馬最後定理》共用前導文件
% 編譯引擎：xelatex

% ── 中文排版 ──────────────────────────────────────────
\usepackage{xeCJK}
\setCJKmainfont[BoldFont={Noto Serif CJK TC SemiBold}]{Noto Serif CJK TC}
\setCJKsansfont{Noto Sans CJK TC}
\setCJKmonofont{Noto Sans Mono CJK TC}
\XeTeXlinebreaklocale "zh"
\XeTeXlinebreakskip = 0pt plus 1pt
% 這些符號改由 CJK 字型出字（拉丁字型缺字）
\xeCJKDeclareCharClass{CJK}{"2500 -> "257F, "2113, "2605, "2606, "2713, "25CF, "25CB, "2194, "2192, "2190, "21D0, "21D2, "21D4, "21A6, "27F9, "27FA, "2208}

% ── 數學 ─────────────────────────────────────────────
\usepackage{amsmath}
\usepackage{amssymb}
\usepackage{amsthm}
\usepackage{mathtools}
\usepackage{tikz-cd}   % 交換圖（第六、七部）

% ── 定理環境（pandoc fenced div 經 theorems.lua 對應到這裡）──
% 單章編譯：定理在章內流水編號。全書編譯時改用 book class 並加 [chapter]。
\theoremstyle{plain}
\newtheorem{thm}{定理}
\newtheorem{prop}[thm]{命題}
\newtheorem{lem}[thm]{引理}
\newtheorem{cor}[thm]{推論}
\newtheorem{conj}[thm]{猜想}
\theoremstyle{definition}
\newtheorem{defn}[thm]{定義}
\newtheorem{exa}[thm]{例}
\theoremstyle{remark}
\newtheorem*{rem}{註}
\renewcommand{\proofname}{證明}

% ── 全書記號巨集（與 notation.md 同步維護）──────────────
\newcommand{\N}{\mathbb{N}}
\newcommand{\Z}{\mathbb{Z}}
\newcommand{\Q}{\mathbb{Q}}
\newcommand{\R}{\mathbb{R}}
\newcommand{\C}{\mathbb{C}}
\newcommand{\F}{\mathbb{F}}                 % 有限體 \F_p
\newcommand{\Qbar}{\overline{\mathbb{Q}}}
\newcommand{\OK}{\mathcal{O}_K}             % 整數環
\newcommand{\Gal}{\operatorname{Gal}}
\newcommand{\GL}{\operatorname{GL}}
\newcommand{\SL}{\operatorname{SL}}
\newcommand{\Frob}{\operatorname{Frob}}
\newcommand{\divides}{\mid}
\newcommand{\notdivides}{\nmid}

% ── 版面 ─────────────────────────────────────────────
\usepackage[a4paper, margin=2.6cm]{geometry}
\usepackage{setspace}
\setstretch{1.18}
\setlength{\parskip}{0.35em}
